%%%%%%%%%%%%%%%%%%%%%%%%%%%%%%%%%%%%%%%%%%%%%%%%
% COPYRIGHT: (C) 2012-2015 FAU FabLab and others
% Bearbeitungen ab 2015-02-20 fallen unter CC-BY-SA 3.0
%%%%%%%%%%%%%%%%%%%%%%%%%%%%%%%%%%%%%%%%%%%%%%%%


\newcommand{\basedir}{fablab-document}
\documentclass[13pt]{\basedir/fablab-document}
\usepackage{tabularx} % Tabellen mit bestimmtem Breitenverhältnis der Spalten
\usepackage{ifthen}
\usepackage{xspace}
\def\tabularnewcol{&\xspace} % hässlicher Workaround von http://tex.stackexchange.com/questions/7590/how-to-programmatically-make-tabular-rows-using-whiledo

\title{Einweisungsliste SLA 3D-Drucker}
\fancyfoot[C]{\small github.com/fau-fablab/sla-drucker-einweisung}
\fancyfoot[L]{Einweisungsliste Nr. \underline{\hspace{3em}}}

\begin{document}
%\maketitle

\textbf{für UniFormation GKtwo mit Phrozen Water-Washable Rapid Black}

Ich bestätige mit meiner Unterschrift verbindlich, dass

\begin{itemize}
\item ich die Regeln und Hinweise, die Betriebsanweisungen für den Drucker
und für das Harz sowie die Anleitung gelesen und verstanden habe
\item ich ein Objekt unter Anleitung erfolgreich ausgedruckt, gewaschen und
nachgehärtet habe
\item die Möglichkeit für Fragen und Diskussion gegeben war
\item die Möglichkeit, eine Kopie der Anleitung zu erhalten, gegeben war
\end{itemize}


% einfach kopiert von Einweisungsliste Lasercutter

\newcounter{i}
\setcounter{i}{1}

\newcommand{\leerezeile}{\hspace{2em} \tabularnewcol \hspace{3em} \tabularnewcol \hspace{2.5em} \tabularnewcol \hspace{2.5em} \tabularnewcol \vbox{\vspace{2em}} \tabularnewcol \tabularnewcol \tabularnewcol \tabularnewline \hline}

\begin{tabularx}{\textwidth}{|l|l|l|l|X|X|X|X|}
  \hline
  \textbf{Nr.} & \textbf{Datum} & \textbf{von} & \textbf{bis} & \textbf{Name} & \textbf{Unterschrift} & \textbf{Einweisender} & \textbf{Unterschrift} \\ \hline
  \whiledo{\value{i}<14}%
  {%
    \stepcounter{i} \leerezeile
  }%
  \leerezeile % doofer Workaround, eigentlich sollte das auch in der Forschleife gehen! Ohne dies wird die Spaltenbegrenzung von Spalte 1 zu weit gezeichnet.
\end{tabularx}

\end{document}
